\documentclass[11pt,a4paper]{article}

\usepackage[T2A]{fontenc}
\usepackage[utf8]{inputenc}
\usepackage[russian,english]{babel}
\usepackage{geometry}
\usepackage{amsmath,amssymb}
\usepackage{hyperref}
\usepackage{booktabs}
\usepackage{enumitem}

\geometry{margin=2.5cm}

\title{Angel-in-Pocket Full Business Agent with Subscriptions\\
A Self-Correcting AI Business Architecture}
\author{qqewq}
\date{\today}

\begin{document}
\selectlanguage{english}
\maketitle

\begin{abstract}
Angel-in-Pocket Full Business Agent with Subscriptions is a backend-first AI business system that transforms a user's idea into a subscription-ready business. The architecture combines multi-agent orchestration, business automation, and a stability layer inspired by GRA-Core to evaluate scenarios, remove inefficient branches, and preserve structural coherence.
\end{abstract}

\noindent\textbf{Keywords:} AI agents, business automation, subscriptions, stability, orchestration, finance, accounting.

\section{Introduction}

The goal of Angel-in-Pocket is to turn a human vision into an operational business. The user provides an idea, preferences, and constraints, while the Angel agent swarm designs the product, business model, subscriptions, and supporting operations.

The system is organized around a stability backend that detects contradictions and removes ineffective branches. This makes the architecture self-correcting rather than purely generative.

\section{Problem Statement}

Many AI systems can propose ideas, but fewer can maintain a business structure over time. The core challenge is to map a human vision into a profitable and stable operating model.

We represent the transformation as:
\[
(V, P, C) \rightarrow B \rightarrow G \rightarrow \Pi,
\]
where $V$ is the vision, $P$ is the preference model, $C$ is the constraint set, $B$ is the business canvas, $G$ is the process graph, and $\Pi$ is the resulting operating policy.

\section{Architecture}

The repository can be viewed as four layers:
\begin{itemize}[leftmargin=1.2cm]
    \item \textbf{gra\_core}: stability evaluation, contradiction detection, scenario ranking.
    \item \textbf{angel}: agents for product, marketing, sales, finance, accounting, and risk.
    \item \textbf{app}: backend services, APIs, and database models.
    \item \textbf{frontend}: user-facing interface and dashboards.
\end{itemize}

Each layer contributes to the overall control loop:
\[
\text{Observe} \rightarrow \text{Evaluate} \rightarrow \text{Nullify} \rightarrow \text{Rebuild}.
\]

\section{Subscriptions Model}

The business model emphasizes recurring revenue. A subscription tier can be written as:
\[
T_i = \{p_i, b_i, r_i, c_i\},
\]
where $p_i$ is price, $b_i$ is billing period, $r_i$ is retention logic, and $c_i$ is churn risk.

The Angel agent evaluates each tier and keeps only those that satisfy profitability and stability constraints.

\section{Finance and Accounting}

To support real-world operation, the system includes finance and accounting scaffolding. The core simplified profit equation is:
\[
\Pi = R - O - \tau,
\]
where $R$ is revenue, $O$ is operating cost, and $\tau$ is taxes.

This layer can later be extended with cashflow forecasting, charts of accounts, ledger postings, and tax-specific logic.

\section{Stability Principle}

The main principle of the architecture is nullification of unstable branches. A process branch should remain only if its utility exceeds its cost:
\[
\text{Keep}(x) \iff U(x) - K(x) > 0.
\]

The stability layer compares alternative branches and selects the one with the strongest structural coherence.

\section{Use Cases}

The system is suitable for:
\begin{itemize}[leftmargin=1.2cm]
    \item subscription-based AI businesses,
    \item productized service systems,
    \item business simulation and planning,
    \item finance-aware agent orchestration,
    \item risk-controlled startup automation.
\end{itemize}

\section{Conclusion}

Angel-in-Pocket is not just a planner or chatbot. It is a self-correcting AI business architecture in which the human provides the vision, the Angel agents execute the business logic, and the stability layer preserves coherence over time.

\bibliographystyle{plain}
\begin{thebibliography}{9}

\bibitem{latex}
LaTeX Project. \emph{LaTeX -- A document preparation system}. \url{https://www.latex-project.org/}

\bibitem{ctan}
CTAN. \emph{The Comprehensive TeX Archive Network}. \url{https://ctan.org/}

\end{thebibliography}

\end{document}